\documentclass[11pt]{article}
\usepackage[a4paper,margin=1in]{geometry}
\usepackage{amsmath,amssymb,bm}
\usepackage{booktabs}
\usepackage{enumitem}
\usepackage{xcolor}
\usepackage[colorlinks=true,linkcolor=blue,citecolor=blue,urlcolor=blue]{hyperref}

\newcommand{\muvec}{\bm{\mu}}
\newcommand{\xfield}{\bm{x}}
\newcommand{\loss}{\mathcal{L}}

\title{Short Literature Review for Time-dependent PPDONet}
\author{Project notes}
\date{\today}

\begin{document}
\maketitle

\section{Goal}

The original PPDONet learns a steady map from disk/planet parameters to final
fields \cite{mao2023ppdonet}:
\begin{equation}
    G_{\mathrm{ss}}:\muvec=(\alpha,h_0,q)\mapsto
    \xfield_{\mathrm{ss}}(r,\theta),
\end{equation}
where \(\xfield\) can mean \(\log\Sigma\), \(v_r\), or \(v_\theta\). The new
project asks for a time-dependent surrogate. The simple version would be
\begin{equation}
    G:\left(\muvec,t\right)\mapsto \xfield(r,\theta,t),
\end{equation}
but the more practical version is a short-time evolution operator
\begin{equation}
    \Phi_{\Delta t}:\left(\muvec,\xfield_n\right)\mapsto \xfield_{n+1}.
\end{equation}
Then long rollout is just a few repeated operator calls:
\begin{equation}
    \xfield_K
    =
    \Phi_{\Delta t_K}\circ\cdots\circ\Phi_{\Delta t_1}
    (\muvec,\xfield_0).
\end{equation}

\paragraph{Main intuition.}
The time-dependent accretion-disk PINN paper \cite{timepinn2025disk} suggests
that long-time training is much easier if we split time into local windows, hard
enforce the initial state inside each window, use periodic angular features, and
rescale outputs. I do not think we should copy it as a PINN directly. The more
useful move is to turn that windowed solver idea into a reusable operator.

\section{Two anchor papers}

\paragraph{PPDONet.}
PPDONet uses the DeepONet idea \cite{lu2021deeponet}: a branch network processes
the system parameters and a trunk network processes coordinates. Their product
gives the field value at \((r,\theta)\). This is nice for disk problems because
the output is a function, not a fixed-resolution image. It already contains some
important engineering choices: log transforms for stiff parameters, periodic
encoding of \(\theta\), and separate models for \(\log\Sigma\), \(v_r\), and
\(v_\theta\).

\paragraph{Time-dependent disk PINN.}
The PINN disk paper is useful mostly as a training recipe. It treats the PDE
residual seriously, but it also makes several very practical moves: train over
short time intervals, enforce the initial condition by construction, use
periodic coordinate features, scale outputs, and balance loss terms. For our
project, these are more valuable than the exact PINN architecture.

\section{Model options}

\begin{table}[h]
\centering
\begin{tabular}{p{0.22\linewidth}p{0.33\linewidth}p{0.33\linewidth}}
\toprule
Model & Why it might work & Main risk \\
\midrule
Time-conditioned DeepONet
& Smallest change from PPDONet. Put local time \(\tau\) in the trunk:
\((r,\sin\theta,\cos\theta,\tau)\).
& It may interpolate time well but still struggle with generalizing across new
input states \(\xfield_n\). \\
\addlinespace
State-conditioned DeepONet
& Add sensors or a downsampled encoding of \(\xfield_n\) to the branch. This
turns the model into a real propagator.
& Sensor choice matters; too few sensors lose structures, too many make training
heavy. \\
\addlinespace
FNO / U-FNO
& Fourier neural operators \cite{li2021fno} and U-FNO-style variants
\cite{wen2022uno} are natural for grid-to-grid time stepping.
& Less resolution-free than DeepONet; polar grid and boundaries need care. \\
\addlinespace
PINO
& Physics-informed neural operator \cite{li2021pino} combines operator learning
with PDE residual losses. Good match if we have sparse time snapshots.
& PDE residuals can dominate or destabilize training without careful weighting. \\
\addlinespace
Multiwavelet operator
& Multiwavelet operators \cite{gupta2021mwt} can represent localized structures
better than pure Fourier bases.
& More implementation work; probably not the first baseline. \\
\bottomrule
\end{tabular}
\caption{Reasonable model swaps for the project.}
\end{table}

\paragraph{Recommended order.}
Start with time-conditioned/state-conditioned DeepONet, because it reuses the
existing code and gives a clean ablation. Then try FNO/U-FNO as the grid-to-grid
baseline. If we later add PDE residuals, call it a PINO-style extension rather
than a totally new model.

\section{How to add time practically}

\subsection{Version 0: time as a coordinate}

The smallest change is
\begin{equation}
    \hat{\xfield}
    =
    N_\theta(\muvec,r,\sin\theta,\cos\theta,\tau).
\end{equation}
This is useful for interpolation between snapshots of the same run, but it is
not yet a proper evolution operator because it does not know the current state.
It is still worth doing as a smoke test.

\subsection{Version 1: hard-IC local window}

For a local window starting from \(\xfield_n\), use
\begin{equation}
    \hat{\xfield}(\muvec,\xfield_n,r,\theta,\tau)
    =
    \xfield_n(r,\theta)
    +
    A(\tau)
    R_\theta(\muvec,E(\xfield_n),r,\sin\theta,\cos\theta,\tau),
\end{equation}
with \(A(0)=0\), for example \(A(\tau)=\tau\). This guarantees
\begin{equation}
    \hat{\xfield}(\tau=0)=\xfield_n.
\end{equation}
This one detail matters a lot. Otherwise the model wastes capacity relearning
the left boundary of every time window.

\subsection{Version 2: few-step rollout}

Train the same local operator on many windows:
\begin{equation}
    \Phi_{\Delta t}(\muvec,\xfield_n)=\hat{\xfield}(\tau=1).
\end{equation}
At inference:
\begin{equation}
    \xfield_{n+1}=\Phi_{\Delta t}(\muvec,\xfield_n).
\end{equation}
This is the most practical ``time-dependent PPDONet'' story: the model is still
an operator, but now the operator maps one field to the next field.

\section{Loss design}

The baseline supervised loss is
\begin{equation}
    \loss_{\mathrm{data}}
    =
    \|\hat{\xfield}(r,\theta,\tau)-\xfield_{\mathrm{FARGO}}(r,\theta,\tau)\|_2^2.
\end{equation}
For time-dependent training, this is not enough. The useful loss stack is:
\begin{equation}
    \loss =
    \lambda_d\loss_{\mathrm{data}}
    +\lambda_{\mathrm{ic}}\loss_{\mathrm{ic}}
    +\lambda_{\mathrm{pde}}\loss_{\mathrm{pde}}
    +\lambda_{\mathrm{roll}}\loss_{\mathrm{roll}}
    +\lambda_{\mathrm{ss}}\loss_{\mathrm{ss}}.
\end{equation}

\paragraph{Initial condition.}
If we use the hard-IC ansatz above, \(\loss_{\mathrm{ic}}\) is basically zero by
construction. This is better than merely adding a penalty.

\paragraph{PDE residual.}
The PINN/PINO idea is to add hydrodynamic residuals using autodiff in
\((r,\theta,\tau)\) \cite{raissi2019pinn,li2021pino}. In practice I would not
turn this on first. First get supervised window prediction working. Then add the
residual with a small weight and monitor whether it improves rollout stability.

\paragraph{Rollout consistency.}
If the model is trained one step at a time, it may still drift in multi-step
rollout. Add a short unrolled loss:
\begin{equation}
    \loss_{\mathrm{roll}}
    =
    \sum_{j=1}^J
    \|\hat{\xfield}_{n+j}-\xfield_{n+j}\|_2^2.
\end{equation}
Use \(J=2\) or \(3\) first. Long unrolls are expensive and can make optimization
annoying.

\paragraph{Steady-state anchor.}
The old PPDONet can act as a soft late-time target:
\begin{equation}
    \loss_{\mathrm{ss}}
    =
    \|\hat{\xfield}(T)-G_{\mathrm{ss}}(\muvec)\|_2^2.
\end{equation}
This is biased if the true simulation has not reached the same steady state, so
it should be weak or only used near late times. Still, it is a nice way to make
the old baseline useful.

\paragraph{Loss balancing.}
PINNs often fail because one term dominates the gradients
\cite{wang2021pinnloss}. The time-dependent disk PINN paper also points in this
direction. For us, a practical recipe is:
\begin{enumerate}[leftmargin=*]
    \item normalize each field separately;
    \item start with data loss only;
    \item add PDE residual after the supervised model is sane;
    \item use gradient-norm or running-loss balancing for \(\lambda\)'s;
    \item use causal/time-window weighting so early-time errors are fixed before
    later-time errors \cite{wang2022causality}.
\end{enumerate}

\section{What looks most promising}

\paragraph{First real experiment.}
Use FARGO snapshots if available. Train a local DeepONet propagator:
\begin{equation}
    (\muvec, E(\xfield_n), r,\sin\theta,\cos\theta,\tau)
    \mapsto
    \xfield(r,\theta,t_n+\tau\Delta t).
\end{equation}
The state encoder \(E\) can initially be just fixed sensors or a coarse grid
flattened through an MLP.

\paragraph{Ablations.}
The important ablations are small and clear:
\begin{enumerate}[leftmargin=*]
    \item no time input vs time input;
    \item soft IC penalty vs hard-IC ansatz;
    \item parameter-only branch vs parameter-plus-state branch;
    \item one-step loss vs two-step rollout loss;
    \item data-only vs data-plus-PDE residual.
\end{enumerate}

\paragraph{My current bet.}
Time as a trunk coordinate will help temporal interpolation, but the real
project lives or dies on the state branch \(E(\xfield_n)\). Without that, the
model is closer to a movie interpolator for a known run. With it, it becomes an
actual evolution operator.

\bibliographystyle{plain}
\bibliography{refs}

\end{document}
